\section{Cloud computing}
\label{sec:cloud_computing}

The most widely accepted definition of cloud computing was proposed by the National Institute of Standards and Technology (NIST). Mell and Grance define cloud computing as ``a model for enabling ubiquitous, convenient, on-demand network access to a shared pool of configurable computing resources (e.g., networks, servers, storage, applications, and services) that can be rapidly provisioned and released with minimal management effort or service provider interaction'' \cite{Mell2011NIST}. This definition has become the de facto standard adopted by both academia and industry.

According to NIST, cloud computing is characterized by five essential characteristics: on-demand self-service, broad network access, resource pooling, rapid elasticity, and measured service \cite{Mell2011NIST}. These characteristics enable organizations to dynamically allocate computing resources according to workload demands while improving scalability and resource utilization. In addition, NIST categorizes cloud deployment into four models: private cloud, public cloud, community cloud, and hybrid cloud. Among these deployment models, hybrid cloud has gained increasing attention because it combines the flexibility and scalability of public cloud services with the control and security offered by private cloud infrastructures.

The widespread adoption of cloud computing has fundamentally transformed modern software development and IT operations. Organizations increasingly deploy native cloud applications using automated development pipelines and distributed infrastructures, making security an essential consideration throughout the software operational lifecycle. As a result, cloud computing serves as the technological foundation for the hybrid cloud environments and SysSecOps process investigated in this research.